\documentclass[12pt]{article}
\usepackage[T1]{fontenc}
\usepackage[utf8]{inputenc}
\usepackage{lmodern}
\usepackage{textcomp}
\usepackage{enumitem}
\usepackage{geometry}
\geometry{margin=1in}
\begin{document}
\begin{center}
Answer Key - History - Graduate Level Assessment 12 \
\small (Answers are ordered exactly as in the question paper.)
\end{center}

\section*{Multiple Choice}
\begin{enumerate}[label=\arabic*.]
\item \textbf{Answer:} A
\\ \textit{Options:} A) simple foragers.; B) maize agriculturalists.; C) conquered by the Aztec.; D) dependent on hunting large game.; E) primarily seafaring traders.
\\ \textit{Note:} The native peoples of the northwest coast of North America were simple foragers because they relied on the abundant natural resources available in their environment, such as fish, shellfish, and wild plants, to sustain their communities without the need for agriculture.
\item \textbf{Answer:} B
\\ \textit{Options:} A) the introduction of Linear A script and the adoption of Linear B script; B) an earthquake and a volcanic eruption; C) the end of the Stone Age and the beginning of the Bronze Age; D) the discovery of bronze on the island, and the replacement of bronze by iron; E) the construction of the Palace of Knossos and its destruction
\\ \textit{Note:} The peak of Minoan civilization is marked by significant geological events, specifically an earthquake and a volcanic eruption, which are believed to have influenced the socio-political dynamics and eventual decline of the civilization.
\item \textbf{Answer:} D
\\ \textit{Options:} A) created intricate pottery and jewelry.; B) discovered the use of fire.; C) invented the wheel.; D) stored food surpluses in good years for distribution in bad years.; E) established trade routes with distant cultures.
\\ \textit{Note:} The correct answer highlights that the ability of Ancestral Pueblo elites to store food surpluses during favorable conditions allowed them to establish control and influence over resource distribution in times of scarcity, thereby solidifying their power and status within the community.
\item \textbf{Answer:} A
\\ \textit{Options:} A) stratigraphic analysis; B) resistivity survey; C) metal detectors; D) ground-penetrating radar; E) underwater archaeology
\\ \textit{Note:} Stratigraphic analysis is the correct answer because it allows archaeologists to examine the layers of soil and artifacts, providing insights into the chronology and extent of large-scale land modifications over time.
\item \textbf{Answer:} A
\\ \textit{Options:} A) bonobos; B) orangutans; C) gibbons; D) chimpanzees; E) baboons
\\ \textit{Note:} Humans share a common ancestor with bonobos, making them our closest living relatives in the primate family tree, which highlights the evolutionary link between our species.
\end{enumerate}

\section*{True / False}
\begin{enumerate}[label=\arabic*.]
\setcounter{enumi}{5}
\item \textbf{Answer:} True
\\ \textit{Options:} A) True; B) False
\\ \textit{Note:} According to the passage: early warning.
\item \textbf{Answer:} False
\\ \textit{Options:} A) True; B) False
\\ \textit{Note:} The correct answer is: 19 th
\item \textbf{Answer:} False
\\ \textit{Options:} A) True; B) False
\\ \textit{Note:} The correct answer is: Ashkenazi Jews
\item \textbf{Answer:} True
\\ \textit{Options:} A) True; B) False
\\ \textit{Note:} According to the passage: 25-year-old design
\item \textbf{Answer:} False
\\ \textit{Options:} A) True; B) False
\\ \textit{Note:} The correct answer is: 27\%
\end{enumerate}

\section*{Long Form Response}
\begin{enumerate}[label=\arabic*.]
\setcounter{enumi}{10}
\item \textbf{Answer:} Via Roma
\\ \textit{Note:} Expected answer: Via Roma
\item \textbf{Answer:} Association of Southeast Asian Nations
\\ \textit{Note:} Expected answer: Association of Southeast Asian Nations
\end{enumerate}

\vfill
\noindent\textit{End of Answer Key}
\end{document}